% !TeX root = ../main.tex
% Fonte: Lezione 7.pdf, pp. 2--9. La costruzione delle potenze reali
% è rinviata nella fonte; la nota ne dà una possibile definizione.
\section{Esponenziali e logaritmi}
\label{sec:exp-log}

\subsection{La funzione esponenziale}

\textbf{Esponenti nulli e negativi.} Per una base $a>0$, le potenze razionali si estendono all'esponente zero ponendo $a^0=1$ e agli esponenti negativi ponendo $a^{-r}=1/a^r$ per $r>0$ razionale.

\textbf{Estensione agli esponenti reali.} La funzione $x\mapsto a^x$ estende a tutti gli esponenti reali le potenze già definite per esponenti razionali.\footnote{Per $a>1$ si può porre $a^x=\sup\{a^r:r\in\Q,\ r\le x\}$; per $0<a<1$ si pone $a^x=(1/a)^{-x}$, mentre $1^x=1$. La verifica completa di questa costruzione e delle proprietà dell'esponenziale reale non è svolta qui: le regole delle potenze, la stretta monotonia e l'immagine $(0,+\infty)$, per $a\ne1$, sono assunte come proprietà note.} Ha dominio $\R$, assume soltanto valori strettamente positivi e passa per $(0,1)$, perché $a^0=1$.
\begin{itemize}
  \item Se $a>1$, è strettamente crescente e biettiva da $\R$ a $(0,+\infty)$.
  \item Se $0<a<1$, è strettamente decrescente e ancora biettiva da $\R$ a $(0,+\infty)$.
  \item Se $a=1$, è la funzione costante $1$ e non è invertibile su $\R$.
\end{itemize}
\textbf{Comportamento ed estremi.} Al crescere dell'argomento, nel caso $a>1$ i valori diventano arbitrariamente grandi; nel caso $0<a<1$ si avvicinano a zero senza mai raggiungerlo. Per argomenti molto negativi i due comportamenti si scambiano. Perciò, quando $a\ne1$, l'infimo è $0$ senza minimo, mentre non esiste un maggiorante reale.

\textbf{Regole delle potenze.} Si ha
\[
a^{x+y}=a^xa^y,\qquad a^{-x}=\frac1{a^x},\qquad a^0=1.
\]
% Raccordo con Lezione 15, pp. 13--15: definizione rigorosa e rappresentazioni nel capitolo sui limiti notevoli.
\textbf{Base $e$.} Una base particolarmente importante è il \emph{numero di Nepero} $e=2{,}718281828\ldots$, definito nella sezione~\ref{sec:numero-nepero}. Con la scrittura \emph{funzione esponenziale}, senza precisare la base, si indica di solito
\[
\exp\colon\R\longrightarrow(0,+\infty),\qquad \exp(x)=e^x,
\]
che è strettamente crescente e biettiva.

\begin{schema}
\begin{tikzpicture}[x=1cm,y=.65cm]
  \draw[asse] (-2.5,0) -- (2.6,0) node[right] {$x$};
  \draw[asse] (0,-.3) -- (0,5.2) node[above] {$y$};
  \draw[intervallo,domain=-2.2:2.2,samples=70] plot (\x,{pow(2,\x)});
  \draw[thick,orange!80!black,domain=-2.2:2.2,samples=70]
    plot (\x,{pow(2,-\x)});
  \fill[coloreschema] (0,1) circle (1.6pt) node[above right] {$1$};
  \node[coloreschema,right] at (1.6,3.1) {$2^x$};
  \node[orange!80!black,left] at (-1.6,3.1) {$(1/2)^x$};
\end{tikzpicture}
\hspace{.6cm}
\begin{tikzpicture}[x=1cm,y=.85cm]
  \draw[asse] (-.3,0) -- (4.8,0) node[right] {$x$};
  \draw[asse] (0,-2) -- (0,2.15) node[above] {$y$};
  \draw[intervallo,domain=.17:4.5,samples=80] plot (\x,{ln(\x)/ln(2)});
  \draw[thick,orange!80!black,domain=.17:4.5,samples=80]
    plot (\x,{-ln(\x)/ln(2)});
  \fill[coloreschema] (1,0) circle (1.6pt) node[below left] {$1$};
  \node[coloreschema,above left] at (4.4,2.1) {$\log_2x$};
  \node[orange!80!black,below left] at (4.4,-1.9) {$\log_{1/2}x$};
\end{tikzpicture}

Esponenziali con basi maggiore e minore di $1$, e rispettivi logaritmi.
\end{schema}

\subsection{Il logaritmo come funzione inversa}

Per $a>0$, $a\ne1$, si definisce il \emph{logaritmo in base $a$} come l'inversa dell'esponenziale:
\[
\log_a\colon(0,+\infty)\longrightarrow\R,\qquad
\log_a x=y\ \Longleftrightarrow\ a^y=x.
\]
\textbf{Dominio, immagine e monotonia.} Il dominio contiene soltanto numeri positivi; l'immagine è tutta $\R$. Il logaritmo è strettamente crescente se $a>1$, strettamente decrescente se $0<a<1$, perché l'inversa conserva il verso della monotonia.

\textbf{Identità del logaritmo.} Per ogni base ammissibile,
\[
\log_a1=0,\qquad \log_a a=1,\qquad
\log_a(a^y)=y\ (y\in\R),\qquad a^{\log_a x}=x\ (x>0).
\]
\textbf{Logaritmo naturale.} Se la base è $e$, si parla di \emph{logaritmo naturale} e si scrive $\ln x=\log_e x$. In particolare $\ln\colon(0,+\infty)\to\R$ è strettamente crescente e inversa di $\exp$.

\section{Funzioni iperboliche}
\label{sec:iperboliche}

\subsection{Definizioni, identità e grafici}
% Fonte: Lezione 7.pdf, pp. 5--7.

Si definiscono \emph{seno iperbolico}, \emph{coseno iperbolico} e \emph{tangente iperbolica} mediante
\[
\sinh x=\frac{e^x-e^{-x}}2,\qquad
\cosh x=\frac{e^x+e^{-x}}2,\qquad
\tanh x=\frac{\sinh x}{\cosh x}.
\]
Sono tutte definite su $\R$: infatti $\cosh x>0$ e il quoziente non ha denominatore nullo. La tangente iperbolica si indica anche con $\operatorname{tgh}$.

\Needspace{6\baselineskip}
\textbf{Proprietà delle tre funzioni.}
\begin{itemize}
  \item \textbf{Seno iperbolico:} dispari, strettamente crescente, con immagine $\R$.
  \item \textbf{Coseno iperbolico:} pari, strettamente decrescente su $(-\infty,0]$ e strettamente crescente su $[0,+\infty)$; ha minimo $1$ nel punto $0$ e immagine $[1,+\infty)$.
  \item \textbf{Tangente iperbolica:} dispari, strettamente crescente, con immagine $(-1,1)$; è limitata, senza massimo né minimo.
\end{itemize}

\begin{proposizione}[Identità iperbolica fondamentale]
Per ogni $t\in\R$ si ha
\[
\boxed{\cosh^2t-\sinh^2t=1}.
\]
\end{proposizione}
\begin{proof}
Sostituendo le definizioni e usando $e^te^{-t}=1$ si ottiene
\[
\cosh^2t-\sinh^2t
=\frac{(e^t+e^{-t})^2-(e^t-e^{-t})^2}{4}
=\frac{4e^te^{-t}}4=1.
\]
\end{proof}

\textbf{Interpretazione geometrica.} Il nome richiama una costruzione geometrica: mentre $P(t)=(\cos t,\sin t)$ percorre la circonferenza $u^2+v^2=1$, il punto $Q(t)=(\cosh t,\sinh t)$ percorre il \emph{ramo destro} dell'iperbole $u^2-v^2=1$, quello con $u\ge1$.

\begin{schema}
\begin{tikzpicture}[scale=1.45]
  \draw[asse] (-1.3,0) -- (1.45,0) node[right] {$u$};
  \draw[asse] (0,-1.25) -- (0,1.4) node[above] {$v$};
  \draw[intervallo] (0,0) circle (1);
  \draw[dotted] (0,0) -- ({cos(40)},{sin(40)});
  \fill[coloreschema] ({cos(40)},{sin(40)}) circle (1pt)
    node[above right] {$P(t)$};
  \node[below] at (0,-1.5) {$u^2+v^2=1$};
\end{tikzpicture}
\hspace{1.3cm}
\begin{tikzpicture}[scale=1.45]
  \draw[asse] (-.3,0) -- (2.25,0) node[right] {$u$};
  \draw[asse] (0,-1.8) -- (0,1.8) node[above] {$v$};
  \draw[intervallo,domain=-1.2:1.2,samples=90]
    plot ({(exp(\x)+exp(-\x))/2},{(exp(\x)-exp(-\x))/2});
  \fill[coloreschema] (1,0) circle (1pt) node[below left] {$1$};
  \fill[coloreschema] ({(exp(.8)+exp(-.8))/2},{(exp(.8)-exp(-.8))/2})
    circle (1pt) node[above left] {$Q(t)$};
  \node[below] at (1,-1.99) {$u^2-v^2=1,\quad u\ge1$};
\end{tikzpicture}
\end{schema}

\begin{proof}[Verifica delle proprietà]
\textbf{Simmetrie.} La parità di $\cosh$ e la disparità di $\sinh$ seguono scambiando $t$ con $-t$.

\textbf{Crescita del seno iperbolico.} Se $u<v$, si ha $e^u<e^v$ ed $e^{-u}>e^{-v}$: dalla differenza delle due quantità segue $\sinh u<\sinh v$.

\textbf{Minimo del coseno iperbolico.} Si ha
\[
\cosh x-1=\frac{(e^{x/2}-e^{-x/2})^2}{2}\ge0,
\]
con uguaglianza soltanto in $x=0$.

\textbf{Monotonia del coseno iperbolico.} Per $0\le u<v$, sviluppando le definizioni si verifica
\[
\cosh v-\cosh u
=2\sinh\frac{u+v}2\,\sinh\frac{v-u}2>0.
\]
La parità dà poi la decrescenza sulla semiretta negativa.

\textbf{Tangente iperbolica.} È utile scrivere $\tanh x=(e^{2x}-1)/(e^{2x}+1)$. Per $0<s<t$,
\[
\frac{t-1}{t+1}-\frac{s-1}{s+1}
=\frac{2(t-s)}{(t+1)(s+1)}>0;
\]
quindi $\tanh$ è strettamente crescente. Il rapporto è sempre compreso fra $-1$ e $1$ e, sostituendo $-x$, cambia segno. Le immagini esatte delle tre funzioni seguono anche dalle formule inverse ricavate qui sotto.
\end{proof}

\begin{schema}
\begin{tikzpicture}[x=.8cm,y=.65cm]
  \draw[asse] (-2.3,0) -- (2.4,0) node[right] {$x$};
  \draw[asse] (0,-3.4) -- (0,3.6) node[above] {$y$};
  \draw[intervallo,domain=-1.9:1.9,samples=80]
    plot (\x,{(exp(\x)-exp(-\x))/2});
  \node[below] at (0,-3.7) {$\sinh x$};
\end{tikzpicture}
\hspace{.3cm}
\begin{tikzpicture}[x=.8cm,y=.65cm]
  \draw[asse] (-2.3,0) -- (2.4,0) node[right] {$x$};
  \draw[asse] (0,-.4) -- (0,3.8) node[above] {$y$};
  \draw[dashed,black!40] (-2.1,1) -- (2.1,1);
  \draw[intervallo,domain=-1.9:1.9,samples=80]
    plot (\x,{(exp(\x)+exp(-\x))/2});
  \node[below right] at (0,1) {$1$};
  \node[below] at (0,-.7) {$\cosh x$};
\end{tikzpicture}
\hspace{.3cm}
\begin{tikzpicture}[x=.8cm,y=1.15cm]
  \draw[asse] (-2.3,0) -- (2.4,0) node[right] {$x$};
  \draw[asse] (0,-1.5) -- (0,1.5) node[above] {$y$};
  \foreach \y in {-1,1} \draw[dashed,black!40] (-2.2,\y) -- (2.2,\y);
  \draw[intervallo,domain=-2.2:2.2,samples=80]
    plot (\x,{(exp(2*\x)-1)/(exp(2*\x)+1)});
  \node[above right] at (0,1) {$1$};
  \node[below right] at (0,-1) {$-1$};
  \node[below] at (0,-1.75) {$\tanh x$};
\end{tikzpicture}
\end{schema}

\subsection{Funzioni iperboliche inverse}
% Fonte: Lezione 7.pdf, pp. 8--9.

\textbf{Settore seno iperbolico.} Il seno iperbolico è biettivo da $\R$ a $\R$. La sua inversa è il \emph{settore seno iperbolico}, indicato con $\operatorname{arsinh}$ o $\operatorname{settsinh}$, e si esprime mediante il logaritmo:
\[
\operatorname{arsinh}x=\ln\!\left(x+\sqrt{x^2+1}\right),
\qquad \operatorname{arsinh}\colon\R\longrightarrow\R.
\]
\begin{proof}[Derivazione]
Risolviamo $x=\sinh y$ e poniamo $t=e^y>0$:
\begin{align*}
x=\frac{e^y-e^{-y}}2
&\ \Longleftrightarrow\ 2x=t-\frac1t\\
&\ \Longleftrightarrow\ t^2-2xt-1=0\\
&\ \Longleftrightarrow\ t=x\pm\sqrt{x^2+1}.
\end{align*}
Poiché $\sqrt{x^2+1}>|x|$, il segno meno darebbe $t<0$ e va scartato; il segno più dà un valore strettamente positivo per ogni reale $x$. Dunque $e^y=x+\sqrt{x^2+1}$ e si applica $\ln$.
\end{proof}

\textbf{Settore coseno iperbolico.} Il coseno iperbolico è invertibile dopo averne ristretto il dominio a $[0,+\infty)$. L'inversa è il \emph{settore coseno iperbolico}:
\[
\operatorname{arcosh}x=\ln\!\left(x+\sqrt{x^2-1}\right),
\qquad \operatorname{arcosh}\colon[1,+\infty)\longrightarrow[0,+\infty).
\]
La notazione alternativa è $\operatorname{settcosh}$.

\begin{proof}[Scelta del ramo]
Da $x=\cosh y$ e $t=e^y$ segue $t^2-2xt+1=0$, con soluzioni $t=x\pm\sqrt{x^2-1}$. Per $x>1$ il segno più dà $t>1$, dunque $y>0$; il segno meno dà $0<t<1$, dunque $y<0$, escluso dalla restrizione scelta. Per $x=1$ le due soluzioni coincidono in $t=1$, cioè $y=0$.
\end{proof}

\textbf{Settore tangente iperbolica.} La tangente iperbolica è biettiva da $\R$ a $(-1,1)$. La sua inversa, il \emph{settore tangente iperbolica}, è
\[
\operatorname{artanh}x=\frac12\ln\!\left(\frac{1+x}{1-x}\right),
\qquad \operatorname{artanh}\colon(-1,1)\longrightarrow\R,
\]
indicata anche con $\operatorname{setttgh}$.

\begin{proof}[Derivazione]
L'equazione $x=(e^{2y}-1)/(e^{2y}+1)$ equivale a $e^{2y}=(1+x)/(1-x)$. Per $-1<x<1$ il rapporto è positivo, quindi si può applicare il logaritmo e dividere per $2$.
\end{proof}

